\section{Evaluation}

\subsection{Evaluation Goals}

The benchmark suite evaluates the routing layer of DIAL under controlled local
conditions. It does not evaluate natural-language answer quality. During these
runs, capability requirements are provided directly by the benchmark and the
selected node returns a dummy response. This keeps the measurement focused on
the distributed system: discovery, profile fetching, scoring, recommendation,
forwarding, and failure handling.

The goal is therefore not to prove that DIAL is ready for a public deployment.
The results show whether the prototype can route requests correctly under
controlled conditions, and where the current discovery and failure-handling
logic starts to weaken. They do not validate real LLM answer quality, public
network behavior, NAT traversal, provider trust, or privacy guarantees.

The latest benchmark snapshot contains 3,810 queries across 38 executed run
folders. Overall, 90.4\% of queries completed successfully, 63.6\% selected the
best active peer according to the pure capability-quality oracle, and 77.8\%
selected a top-3 peer. The average selected-to-oracle quality ratio was
81.8\%.
The failures were concrete: 276 queries ended with \texttt{no\_suitable\_node}
and 91 ended with a timeout.

The main metrics are:

\begin{itemize}
  \item \textbf{success rate}: fraction of queries returning a valid response;
  \item \textbf{top-1 accuracy}: whether the selected peer was the best active
  peer by pure capability quality;
  \item \textbf{top-3 accuracy}: whether the selected peer was among the three
  best active peers by pure capability quality;
  \item \textbf{quality ratio}: selected peer quality divided by oracle peer
  quality;
  \item \textbf{routing latency}: time spent in discovery, profile exchange,
  scoring, and forwarding decisions;
  \item \textbf{p95 end-to-end latency}: tail latency observed by the benchmark
  client;
  \item \textbf{candidate count}: how many peers were evaluated by the router;
  \item \textbf{candidate coverage}: evaluated candidates divided by active
  nodes.
\end{itemize}

For readability, rates and quality ratios are written as percentages in the
text. The tables and plots keep the same quantities on a 0--1 scale.

Top-1 accuracy is deliberately strict. DIAL is not meant to send every request
to the same highest-quality peer; it should choose a suitable peer while also
considering reachability, latency, freshness, and eventually load. For this
reason, quality ratio, top-3 accuracy, and candidate coverage are often more
useful than top-1 alone.

One detail matters when reading the latency numbers. The benchmark runner uses
a client-side timeout configured per scenario, while each node also has its own
peer-query timeout. When a p95 value is close to 90, 120, 150, or 180 seconds,
it usually means some requests waited until a benchmark timeout. These values
should therefore be read as timeout pressure, not as normal response time for a
successful route.

\subsection{Clean and Latency Runs}

The first result is positive: the basic routing path works. In clean 8-node
runs, every query completed and the selected peer was almost always the oracle
peer. Moderate simulated latency also remained stable. The 50--150 ms latency
run still achieved 100\% success and a 99.6\% quality ratio.

\begin{table}[H]
\centering
\caption{8-node baseline and latency results.}
\begin{tabular}{lrrrrrr}
\toprule
\textbf{Scenario} & \textbf{Queries} & \textbf{Success} & \textbf{Top-1} & \textbf{Top-3} & \textbf{Quality} & \textbf{P95 ms} \\
\midrule
No latency & 120 & 1.000 & 0.967 & 1.000 & 0.991 & 277 \\
50--150 ms latency & 120 & 1.000 & 0.983 & 1.000 & 0.996 & 4,262 \\
100--500 ms latency & 120 & 0.917 & 0.867 & 0.917 & 0.899 & 120,023 \\
\bottomrule
\end{tabular}
\end{table}

High latency is different. At 100--500 ms simulated delay, success fell to
91.7\%, and the p95 end-to-end latency reached the scenario's 120-second
timeout range. This does not mean the score formula suddenly became bad, or
that successful answers normally take that long. It means too many requests
wait on slow network operations or failed forwards until the benchmark timeout
is reached. Better timeout handling and more parallel discovery would matter
more than small score-weight changes here.

In short, DIAL handles clean and moderate-latency routing well, but it should
fail faster when network operations become very slow.

\subsection{Scalability and Candidate Visibility}

The scalability campaign is the most important stress test. It keeps latency at
50--150 ms and grows the network from 8 to 128 nodes. Success remains high, but
quality degrades as the network grows: at 128 nodes, success is still 84.2\%,
but the quality ratio is 59.1\% and top-1 accuracy is only 17.5\%.

\begin{figure}[H]
\centering
\begin{tikzpicture}
\begin{axis}[
  width=0.92\textwidth,
  height=6.1cm,
  xlabel={Number of nodes},
  ylabel={Metric value},
  xmin=0,
  xmax=136,
  ymin=0,
  ymax=1.05,
  grid=major,
  legend style={
    at={(0.5,1.03)},
    anchor=south,
    draw=none,
    fill=none,
    legend columns=2,
    /tikz/every even column/.append style={column sep=0.45cm}
  }
]
\addplot[color=deepblue,mark=*] coordinates {
  (8,1.000) (16,0.950) (32,0.942) (64,0.858) (128,0.842)
};
\addplot[color=epflred,mark=square*] coordinates {
  (8,0.989) (16,0.924) (32,0.804) (64,0.618) (128,0.591)
};
\addplot[color=grayblack,mark=triangle*] coordinates {
  (8,0.958) (16,0.708) (32,0.433) (64,0.225) (128,0.175)
};
\addplot[color=blue!55!cyan,mark=diamond*] coordinates {
  (8,0.353) (16,0.277) (32,0.187) (64,0.083) (128,0.038)
};
\legend{Success, Quality ratio, Top-1, Candidate coverage}
\end{axis}
\end{tikzpicture}
\caption{Scaling exposes the candidate-visibility bottleneck. The router still often completes the query, but it evaluates a shrinking share of the active network as node count grows.}
\end{figure}

The reason is visible in candidate coverage. At 8 nodes, the router evaluates
about 35\% of active nodes. At 128 nodes, it evaluates only about 3.8\%. Since
the oracle is computed over all active profiles, larger networks create more
opportunities for a better node to exist somewhere outside the evaluated set.
This should not be read as a failure by itself. In a larger decentralized
network, finding the exact top-1 or top-3 peer naturally becomes harder because
there are more possible competitors, and DIAL should not concentrate all
traffic on the same strongest nodes. The goal is to find a suitable peer and
eventually spread load across the network. Still, the simultaneous drop in
candidate coverage and quality ratio shows that discovery should improve for
larger deployments. The implemented recommendation step already moves in this
direction by letting peers suggest additional candidates, but it remains a
conservative one-hop mechanism with a small per-peer limit. Further work should
tune DHT/provider lookup behavior, return or refresh more useful providers per
capability, improve recommendation diversity, and explore bounded multi-hop
candidate expansion.

\subsection{Dropout}

The dropout campaign starts a 24-node network, kills a fraction of non-entry
nodes, waits briefly, and then sends queries. Dropout reduces the candidate
pool and increases \texttt{no\_suitable\_node}. The trend is not perfectly
monotonic, because the removed nodes depend on the sampled dropout set, but the
direction is clear: fewer active useful peers lead to weaker quality and long
tail latency.

\begin{figure}[H]
\centering
\begin{tikzpicture}
\begin{axis}[
  width=0.92\textwidth,
  height=6.1cm,
  xlabel={Dropout rate},
  ylabel={Rate or ratio},
  ymin=0,
  ymax=1.05,
  grid=major,
  symbolic x coords={0,25,50,75},
  xtick=data,
  legend style={
    at={(0.5,1.03)},
    anchor=south,
    draw=none,
    fill=none,
    legend columns=3,
    /tikz/every even column/.append style={column sep=0.45cm}
  }
]
\addplot[color=deepblue,mark=*] coordinates {(0,0.958) (25,0.892) (50,0.833) (75,0.875)};
\addplot[color=epflred,mark=square*] coordinates {(0,0.904) (25,0.801) (50,0.747) (75,0.716)};
\addplot[color=grayblack,mark=triangle*] coordinates {(0,0.858) (25,0.775) (50,0.783) (75,0.842)};
\legend{Success, Quality ratio, Top-3}
\end{axis}
\end{tikzpicture}
\caption{Dropout lowers routing quality and keeps tail latency near timeout-bound values. The 75\% case has a small success rebound in this run, but its quality ratio and candidate count remain weak.}
\end{figure}

At 50\% dropout, success is 83.3\%, quality ratio is 74.7\%, and p95 end-to-end
latency is about 60 seconds. The average candidate count falls from 5.63 with
no dropout to 2.71 at 50\% dropout and 1.86 at 75\% dropout. This is still
encouraging: even after many peers disappear, the prototype usually returns an
answer. The cost is weaker candidate visibility, which makes weaker selections
more likely. Dropout should mainly be read as a discovery and liveness problem.

\subsection{Mixed-Capability Queries}

Mixed-capability queries are important because real prompts often require more
than one skill. The clean 24-node mixed run is strong: all 45 queries complete,
top-3 accuracy is 100\%, and the quality ratio is 99.0\%. This shows that the
weighted capability model can work when the network is visible and fast.

\begin{table}[H]
\centering
\caption{Mixed-capability results under clean, latency, and dropout settings.}
\begin{tabular}{lrrrrr}
\toprule
\textbf{Scenario} & \textbf{Success} & \textbf{Top-1} & \textbf{Top-3} & \textbf{Quality} & \textbf{P95 ms} \\
\midrule
24 nodes, clean & 1.000 & 0.956 & 1.000 & 0.990 & 2,325 \\
16 nodes, 50--150 ms & 0.756 & 0.600 & 0.622 & 0.687 & 90,029 \\
16 nodes, 100--500 ms & 0.711 & 0.489 & 0.533 & 0.618 & 120,026 \\
24 nodes, 25\% dropout & 0.733 & 0.511 & 0.600 & 0.662 & 120,026 \\
\bottomrule
\end{tabular}
\end{table}

The same workload becomes much harder once latency or dropout is added. Mixed
queries force the router to find peers that are good across several capability
dimensions, so narrow candidate discovery hurts more than in single-capability
queries. The practical improvement is to collect candidates per relevant
capability, keep more than one provider per capability, and then score the
combined requirement after deduplication. The rows with p95 values near 120
seconds should again be read as timeout-bound failures, not as normal response
times for completed answers.

In short, the clean run validates the weighted capability representation. The
stressed runs show that multi-capability routing needs broader candidate
collection before the scoring formula can really help.

\subsection{Policy Comparison}

The policy runs compare five routing policies in four contexts: clean,
high-latency, 50\% dropout, and mixed-capability workloads. The winner changes
by context, which is useful: it shows that policy design matters. But it also
shows that policy tuning alone cannot fix low candidate coverage or long
timeouts.

\begin{table}[H]
\centering
\caption{Best policy by context in the latest benchmark snapshot.}
\begin{tabular}{llrrrr}
\toprule
\textbf{Context} & \textbf{Best policy} & \textbf{Success} & \textbf{Top-1} & \textbf{Quality} & \textbf{P95 ms} \\
\midrule
Clean & \texttt{latency\_aware} & 1.000 & 0.842 & 0.948 & 1,365 \\
High latency & \texttt{capability\_only} & 0.925 & 0.517 & 0.813 & 60,039 \\
50\% dropout & \texttt{balanced} & 0.858 & 0.608 & 0.777 & 66,622 \\
Mixed & \texttt{reliability\_aware} & 0.867 & 0.689 & 0.787 & 150,014 \\
\bottomrule
\end{tabular}
\end{table}

In clean runs, all policies reach 100\% success and differences are small. Under
high latency, \texttt{capability\_only} performs best in this dataset, which is
slightly counterintuitive but informative: adding latency penalties does not
create better candidates, and under slow network operations it can still spend
time evaluating or forwarding to poor paths. Under dropout, \texttt{balanced}
does best, while mixed workloads benefit most from \texttt{reliability\_aware}.
These results should not be read as a final policy choice. They are better read
as evidence that future policy studies should use more contrasted weights and
scenarios designed to isolate one trade-off at a time. The large p95 value in
the mixed row is also tied to the benchmark's 150-second scenario timeout.
The comparison is still useful: policy tuning can improve selection among
visible peers, but it cannot compensate for candidates that were never
discovered.

\subsection{Evaluation Summary}

The benchmark supports four main conclusions:

\begin{enumerate}
  \item The end-to-end DIAL routing path works. Clean and moderate-latency
  results show that nodes can advertise capabilities, discover peers, select a
  suitable node, and return a response.
  \item The main weakness is candidate visibility. As the network grows, the
  router evaluates a smaller share of active nodes, so the best peer is more
  likely to remain unseen.
  \item Dropout and high latency mostly hurt through stale or slow paths. They
  increase failed routes, timeout-bound latency, and quality loss.
  \item Mixed-capability routing works in clean conditions but needs broader
  candidate collection when latency, dropout, or larger networks are involved.
\end{enumerate}

The report therefore focuses on the plots and tables that explain these
points. The short version is that DIAL routes well when the network is small
and visible; the next engineering work is broader discovery through DHT tuning
and bounded recommendations, faster failed-peer handling, and more deliberate
multi-capability candidate collection.
